\section{The Spinglass model}
We now move on a more complex model, in order to approach the structure of the Hopfield network.
\subsection{The Edwards–Anderson model}
The Edwards–Anderson (EA) model -- first introduced by Edwards and Anderson in 1975 \cite{SFEdwards_1975} -- is a finite dimensional spin glass model with short range interaction. It is not a mean field spin glass model with infinite range interactions such as the Sherrington–Kirkpatrick model \cite{SK_model}. The model is formally described as follows:
\begin{equation}
  H_\text{EA}=\sum_{\langle i,j \rangle}J_{ij}S_iS_j.
  \label{eq:EA_model}
\end{equation}
Where $J_{ij}$ is a random variable taking its values in $\{-1,\,1\}$ with equal probability. Thus, it can easily be represented by a \Bits.
Under a spin flip ($S_k \to - S_k$), the energy shift is:
\begin{equation}
  \Delta E_k=2 S_k \sum_{i \in \partial k}J_{ki}S_i.
  \label{eq:DeltaE_EA}
\end{equation}
The structure being really similar to that of \eqref{eq:DeltaE_ising}, we can perform the same derivation as for the Ising model and find the spin flip condition:
\begin{equation}
  \lfloor \rho \rfloor \geq S_k  \sum_{i \in \partial k}J_{ki}S_i,\\
  \label{eq:flip_EA_int}
\end{equation}
with this time $\rho \sim \mathcal{E}(2\beta)$.
\subsection{Bitwise Implementation}
We now need to express the product $J_{ik}S_kS_i$ in a binary way. As shown in the previous section, the product $S_kS_i$ can be re-written $(-1+2B_k)(-1+2B_i)=-1+2 C_{ki}$ with $C_{ki}= B_k \odot B_i$. Thus, by writing $J_{ki}=-1+2K_{ki}$ with $K_{ki}$ a \Bits, we have:
\begin{equation*}
  \begin{aligned}
  J_{ik}S_kS_i =& (-1+2K_{ki})(-1+2 C_{ki})\\
               =& -1 + 2D_{ki},\\
  \end{aligned}
\end{equation*}
with $D_{ki}=K_{ki}\odot C_{ki}= K_{ki} \odot B_k \odot B_i =  K_{ki} \oplus B_k \oplus B_i$. As in the previous section, the flip condition is then:
\begin{equation}
  \lfloor \rho \rfloor + n_k \geq 2 \sum_{i \in \partial k}D_{ki}.
  \label{eq:flip_SG_bitwise}
\end{equation}
Both the bitwise and the classic implementation are thus similar to this of the Ising model, and we will not present them further in this report. The bitwise algorithm for the spin glass model can be found in the appendix, see \Cref{alg:SG_bitwise}.\\
The bitwise implementation is not validated here, as the bitwise framework was already tested in the previous section, and the main interest of this model lies in bridging the complexity gap between the Ising model and the Hopfield network. We nonetheless verified that the standard and the bitwise implementations produce identical spin configurations when initialized from the same initial state and driven by the same sequence of random numbers.
\begin{algorithm}[t]
\DontPrintSemicolon
\caption{Bitwise spin glass Metropolis}

\KwIn{Bit-replicated spins $\{B_i\}$, couplings $\texttt{Couplings}$, neighbor lists $\partial i$, inverse temperature $\beta$}
\KwOut{Updated spin configurations $\{B_i\}$}

\BlankLine
\textbf{Preallocated variables:} $\mathtt{Bits}\ \texttt{xnor},\ \texttt{mask}$\;
\phantom{\textbf{Preallocated variables:}} $\mathtt{UnsignedInt}\ T_i,\ \texttt{threshold}$\;

\BlankLine

\For{$\mathrm{sweep}=1$ \KwTo $N_{\mathrm{sweeps}}$}{
    \For{$\mathrm{step}=1$ \KwTo $N$}{
        Draw $k$ among $N$\;
        Draw $\lfloor \rho \rfloor$ with $\rho \sim \mathcal{E}(2\beta)$

        \eIf{$\lfloor \rho \rfloor \ge n_k$}{
            $B_k \leftarrow \neg B_k$ \tcp*{escape condition: unconditional flip across all replicas}
        }{
            $T_k \leftarrow \mathbf{0}$

            $\mathtt{threshold} \leftarrow \mathbf{0}$

            \For{$i \in \partial k$}{
                $\mathtt{xnor}\ \leftarrow K_{ki} \oplus B_k \oplus B_i$
                \tcp*{bitwise XNOR across replicas}

                $T_k \mathrel{+}= \mathtt{xnor}$
                \tcp*{increment agreement count per replica}
            }

            $T_k \leftarrow 2T_k$

            $\mathtt{threshold} \leftarrow \lfloor \rho \rfloor +n_k$

            $\mathtt{mask} \leftarrow (T_k \le \mathtt{threshold})$
            \tcp*{per-replica flip condition}

        $B_k \leftarrow B_k \oplus \mathtt{mask}$ \tcp*{bitwise spin flip (XOR operation): only replicas}
        \tcp*[r]{for which $\mathtt{mask}^{(\ell)}=1$ are flipped}
        }
    }
}

\label{alg:SG_bitwise}
\end{algorithm}
\FloatBarrier
\subsection{Performance test}
 We now test the performance of the implementation, for several lattice sizes and the temperatures, this time with $N_\text{sweeps}=2^{14}$. The results are shown in \Cref{fig:speedup_spinglass}. We notice the same evolution of the normalized time for both implementation as the ones observed for the Ising model. The shape of the speedup factor is similar to this of the Ising model, but its evolution is steadier, with what appears to be a saturation for the highest temperatures. The time per sweep also remains proportional to the system size, the complexities of the implementations thus should be similar to those of the Ising model.

\clearpage
\newpage
\FloatBarrier


\nolinenumbers
\begin{figure}[htbp]
\centering
% ── Top row ───────────────────────────────────────────────────────────────
\begin{minipage}[b]{0.49\textwidth}
    \centering
    \includegraphics[width=\columnwidth]{speedup_curves_spinglass.pdf}
    \subcaption{Speedup factor evolution.}
    \label{fig:speedup_curve_spinglass}
\end{minipage}
\hfill
\begin{minipage}[b]{0.49\textwidth}
    \centering
    \includegraphics[width=\columnwidth]{time_normalized_spinglass.pdf}
    \subcaption{Normalized time evolution.}
    \label{fig:normalised_time_spinglass}
\end{minipage}

\vspace{0.5em}

% ── Bottom row ────────────────────────────────────────────────────────────
\begin{minipage}[b]{\textwidth}
    \centering
    \includegraphics[width=\columnwidth]{time_total_vs_N_spinglass.pdf}
    \subcaption{Sweep time depending on the lattice size for all temperatures.}
    \label{fig:sweep_time_vs_N_spinglass}
\end{minipage}
\caption{Performance analysis of the bitwise and classical Spin Glass Metropolis implementations.}
\label{fig:speedup_spinglass}
\end{figure}
\linenumbers
\FloatBarrier
